\subsubsection{HEART-FID Illustration}\label{sec:heartfid_results}

Using the HEART-FID trial as preliminary data, the independence-based FORSS design yielded a fixed sample size of $m=n=1244$ per arm (total $N=2488$). Under independence, the average simulated win ratio was 1.149, with calculated and empirical WR powers of 85.00\% and 84.43\%, respectively. 
When the censoring-aware observed association targets $(-0.2211, 0.5173, -0.1032)$ were used directly as working latent correlations on the raw generator scale, the observed simulated associations were (-0.15, 0.55, -0.06), the average simulated WR was 1.132, with corresponding calculated and empirical WR powers of 76.20\% and 76.59\%. 
By contrast, the calibration-based Gaussian copula approach required a substantially stronger latent correlation matrix $(-0.30, 0.49, -0.17)$ to reproduce the observed Harrell-C/Kendall targets more closely, yielding simulated observed associations (-0.21, 0.53, -0.10). Under this calibrated dependence structure, the average simulated WR was 1.130, with calculated and empirical WR powers of 74.87\% and 75.13\%. 
Across all three scenarios, the empirical WR type I error remained close to the nominal 5\% level (4.85\%, 5.28\%, and 5.09\% for independence, direct input, and calibration, respectively). These results indicate that the direct-input analysis and the calibration-based latent-copula analysis target different correlation quantities and therefore need not yield the same power profile.
